\documentclass[11pt]{extarticle}
\usepackage{graphicx} % Required for inserting images
\usepackage{microtype}
\usepackage{float}
\usepackage{amsmath}
\usepackage{amsthm}
\usepackage{amsfonts}
\usepackage{amssymb}
\usepackage{enumerate}
\usepackage{enumitem}
\usepackage{fancyhdr}
\usepackage{graphicx}
\usepackage{mdframed}
\usepackage{multicol}
\usepackage{mathtools}
\usepackage{listings}
\usepackage{verbatim}
\usepackage{texdraw}
\usepackage{tikz}
\usepackage{titling} 
\usepackage{bm}
\usepackage{booktabs}
\usepackage{hyperref}
\usepackage{indentfirst}
\usepackage{caption}
\usepackage{titlesec}
\usepackage{xcolor}
\definecolor{navyblue}{RGB}{0,32,91}
\hypersetup{colorlinks=true, linkcolor=navyblue, urlcolor=navyblue, citecolor=navyblue}
\usepackage[letterpaper, margin=0.65in]{geometry}
\linespread{1.0}

% Reduce vertical space around floats
\setlength{\textfloatsep}{6pt plus 2pt minus 4pt}
\setlength{\floatsep}{4pt plus 2pt minus 2pt}
\setlength{\intextsep}{6pt plus 2pt minus 2pt}

% Reduce caption spacing
\captionsetup{skip=3pt, font=small, labelfont=bf}

% Compact section/subsection heading spacing + navy color
\titleformat{\section}{\large\bfseries\color{navyblue}}{\thesection}{1em}{}[\vspace{1pt}\color{navyblue}\titlerule]
\titleformat{\subsection}{\normalsize\bfseries\color{navyblue}}{\thesubsection}{1em}{}
\titlespacing*{\section}{0pt}{8pt plus 2pt minus 2pt}{4pt plus 1pt}
\titlespacing*{\subsection}{0pt}{6pt plus 2pt minus 2pt}{3pt plus 1pt}

% Header / footer setup
\pagestyle{fancy}
\fancyhf{}
\fancyhead[L]{\small\itshape Machine Learning in NBA Over/Under Betting}
\fancyhead[R]{\small Page \thepage}
\fancyfoot[C]{\small Business 35137 $|$ Chicago Booth $|$ Winter 2026}
\renewcommand{\headrulewidth}{0.4pt}
\renewcommand{\footrulewidth}{0.4pt}


\newcommand{\NN}{\mathbb{N}}
\newcommand{\ZZ}{\mathbb{Z}}
\newcommand{\RR}{\mathbb{R}}
\newcommand{\CC}{\mathbb{C}}
\newcommand{\QQ}{\mathbb{Q}}
\newcommand{\FF}{\mathbb{F}}
\newcommand{\PP}{\mathbb{P}}
\newcommand{\one}{\mathbb{1}}
\newcommand{\LL}{\mathcal{L}}
\newcommand{\JJ}{\mathcal{J}}
\newcommand\norm[1]{\lVert#1\rVert}
\newcommand\set[1]{\left\lbrace #1 \right\rbrace}
\newcommand{\dist}{\operatorname{dist}}
\newcommand{\nt}{\operatorname{int}}
\newcommand{\mean}{\operatorname{Mean}}
\newcommand{\vol}{\operatorname{vol}}
\newcommand{\dpartial}[2]{\frac{\partial #1}{\partial #2}}
\newcommand\independent{\protect\mathpalette{\protect\independenT}{\perp}}
\newcommand{\indep}{\perp \!\!\! \perp}
\newcommand{\LoadEPS}[4]{%
\setbox0=\hbox{\includegraphics{#1.eps}}
\LoadPDF{#1.pdf}{#2}{#3}{#4}}
\newcommand{\LoadPDF}[4]{%
\pdfximage {#1}
\setbox0=\hbox{\pdfrefximage\pdflastximage} 
\dimen0=#2\wd0         
\dimen1=#2\ht0
\pdfximage              
    width \dimen0 {#1}  
\divide\dimen0 by 2
\divide\dimen1 by 2
\dimen2=#3 \divide\dimen2 by 2
\dimen3=#4 \divide\dimen3 by 2
\vbox to #4{\vskip\dimen3\hbox to#3{\hskip\dimen2%
\hbox to0pt{\vtop to0pt{\vskip-\dimen1 %
\hbox to 0pt{\hskip -\dimen0\pdfrefximage \pdflastximage\hss}\vss}\hss}%
\hss}\vss}%
}

\newenvironment{problem}[2][Problem]{\begin{trivlist}
\item[\hskip \labelsep {\bfseries #1}\hskip \labelsep {\bfseries #2.}]}{\end{trivlist}}
\newenvironment{solution}{\emph{\textbf{Solution}:}}{\qed}  
\newenvironment{proof}{\emph{\textbf{Proof}:}}{\qed}  

\title{Final Project Paper}
\author{Chris Mulligan}
\date{March 2026}

\begin{document}

\begin{titlepage}
    \centering
    \vspace*{\fill}
    \small

    % Title of the Paper
    {\LARGE \textbf{Machine Learning in NBA Over/Under Betting}} \\[2cm]

    % Authors
    {
        Alan Donnelly \\
        Master's in Financial Mathematics \\[0.5cm]
        
        Andrew McLaughlin \\
        Master's in Financial Mathematics \\[0.5cm]
        
        Robert Asgeirsson \\
        Master's in Financial Mathematics \\[0.5cm]
        
        Chris Mulligan \\
        Master's in Financial Mathematics
    } \\[2cm]

    \includegraphics[width=0.5\textwidth]{University Logo_1Color_Maroon_RGB.png} \\[0.5cm]

    % Class and University
    {\large Business 35137: Machine Learning in Finance} \\[0.5cm]
    {\large The University of Chicago Booth School of Business} \\[0.5cm]

    % Date
    {\large Winter 2026}
    
    \vspace*{\fill}
\end{titlepage}

\newpage
\tableofcontents

\newpage

\section{Introduction}
In sports betting, a total (also called an over/under) is the sportsbook's prediction of the combined points scored by
both teams in a game. Bettors wager on whether the actual combined score will be over or under this number.
Standard odds of -110 mean a bettor must risk \$110 to win \$100, implying the sportsbook takes a commission (the
vig) on every bet. To break even at -110 odds, a bettor must win at least 52.4\% of bets.
The Vegas line is set by professional oddsmakers who incorporate team statistics, injuries, matchup history, public
betting patterns, and other market signals. It is widely regarded as the most efficient public forecast of game
outcomes. Any model that claims to beat the line must find a small, repeatable edge that the market has overlooked. \\
This report investigates whether machine learning models can predict NBA game totals more accurately than the
Vegas line, and under what conditions selective betting strategies can generate positive returns. We train six model
classes spanning linear, dimension-reduced, tree-based, and kernel-based approaches using walk-forward validation
on 11 NBA seasons (2013/14 through 2023/24), then test the nominated strategy on a held-out 2024/25 season that
was never used in any modeling decision.

\subsection{Approach}
We construct lagged, within-season team features from box score statistics to avoid lookahead bias. Six model
classes are trained: Ridge, Lasso, ElasticNet, Principal Component Regression (PCR), XGBoost, and RBF+Ridge
(Radial Basis Function kernel approximation combined with Ridge regression). Models are evaluated via
walk-forward backtesting with strict temporal ordering, tuned on betting accuracy rather than regression error, and
subjected to rigorous statistical testing including Bonferroni-corrected significance tests and McNemar's paired
comparison. Selective betting with confidence thresholds is explored to concentrate bets where the model's edge is
strongest. Finally, bankroll simulations with 10 sizing rules stress-test whether statistical edge translates into
economic viability, and real-world frictions (slippage, market depth, account restrictions) are quantified.

\section{Data}
Two primary data sources cover 18 NBA seasons (2007/08 through 2024/25). Betting data contains approximately
23,000 games with the Vegas total line, point spread, and final score. NBA box score data provides team-level
statistics for each game: points, field goal percentage, three-point percentage, free throw percentage, rebounds,
assists, steals, blocks, turnovers, and personal fouls. Historical box scores cover 2007/08 through 2022/23; the
2023/24 and 2024/25 seasons are retrieved via the NBA Stats API and merged with the archive. \\
The two sources are merged on date and team matchup. Team abbreviations are standardized to account for
franchise relocations (e.g., NJN to BKN, SEA to OKC). Each game is assigned a season-end year. The 2024/25
season (1,293 games) is held out entirely from model training and evaluation; the remaining 21,397 games form the
walk-forward training and test pool.
\begin{figure}[H]
    \centering
    \includegraphics[width=1\linewidth]{figure 1.png}
    \caption{Data coverage by season, NBA scoring trend, and Vegas line accuracy}
    \label{fig: 1}
\end{figure}


\section{Feature Engineering}
Feature engineering transforms raw game data into 114 numerical inputs organized into seven categories: Vegas line features (3), totals and scoring (18), shooting efficiency (24), rebounding and hustle (16), playmaking (8), schedule
and rest (11), context indicators (6), and interaction features (28). All features are strictly lagged: rolling means are
computed over each team's previous games within the same season using a shift(1) operation, ensuring no future
information leaks into any prediction. \\
Rolling averages are computed at two window sizes: L5 (last 5 games) to capture short-term form, and L10 (last 10
games) for medium-term stability. For the first game of each team's season, no prior data is available, so these
games are dropped. An early-season indicator variable flags the first 10 games of each team's season, where rolling
statistics are based on limited observations and are therefore noisier. \\
\begin{figure}[H]
    \centering
    \includegraphics[width=1\linewidth]{figure 2.png}
    \caption{Feature composition, target distribution, and Vegas line vs. actual outcomes}
    \label{fig: 2}
\end{figure}

Univariate correlation analysis confirms that Vegas-derived features and recent scoring averages are the strongest
individual predictors ($r = 0.55-0.65$), while schedule and rest features show near-zero marginal correlation. Interaction
features (e.g., pace x pace, spread x total) are included to help linear models capture non-additive effects. \\
\begin{figure}[H]
    \centering
    \includegraphics[width=1\linewidth]{figure 3.png}
    \caption{Top and bottom features ranked by absolute correlation with actual game totals}
    \label{fig: 3}
\end{figure}

\section{Model Training and Walk-Forward Evaluation}
Walk-forward backtesting enforces strict temporal ordering: for each of 11 test seasons (2013/14 through 2023/24),
the model trains on all available data before that season and predicts the test season. The training set grows by
approximately 1,300 games per iteration. Hyperparameters are tuned via 5-fold TimeSeriesSplit within the training
set, optimizing betting accuracy rather than RMSE, since the downstream task is predicting over/under outcomes. \\
A key design choice is exponential recency weighting controlled by a half-life parameter, tuned jointly with model
hyperparameters. Each half-life/hyperparameter combination is evaluated during cross-validation, and the
combination that maximizes betting accuracy is selected for that season. This allows the models to adapt to the
NBA's evolving pace and scoring trends without manually specifying a recency window. \\
Pooling predictions across all 11 test seasons yields approximately 11,000-13,000 out-of-sample bets per model. A no-bet filter excludes games where the model's prediction is within 0.5 points of the Vegas line. \\
\begin{figure}[H]
    \centering
    \includegraphics[width=1\linewidth]{figure 4.png}
    \caption{Walk-forward results: betting accuracy, ROI, and overfitting check (train vs OOS RMSE) for all six models}
    \label{fig: 4}
\end{figure}

\subsection{Key Findings}
All six models cluster between $50.7\%$ and $52.0\%$ accuracy. The linear models (Ridge at $51.7\%$, Lasso and ElasticNet
at $52.0\%$) achieve the highest accuracies, while XGBoost ($50.7\%$) and RBF+Ridge ($50.9\%$) trail behind. No model
reaches the $52.4\%$ break-even threshold when betting every qualifying game, so all ROI values are negative. This is
the expected result in an efficient market. \\
Simple models outperform complex ones. The overfitting check (Figure \ref{fig: 4}, Panel 3) reveals that linear models show
minimal train-vs-OOS RMSE gaps ($\sim0.4$ points), while XGBoost shows a 2.3-point gap, confirming it memorizes
training patterns that do not generalize. In an efficient market where exploitable signal is small and noisy, the
additional flexibility of complex models becomes a liability. \\
\begin{figure}[H]
    \centering
    \includegraphics[width=1\linewidth]{figure 5.png}
    \caption{Per-season accuracy and consistency across 11 walk-forward test seasons.}
    \label{fig: 5}
\end{figure}

Per-season analysis reveals that accuracy is volatile (48-55\% range), with no model consistently beating $52.4\%$
across all seasons. The linear models tend to move together, suggesting they extract similar signal. Lasso achieves
the highest season-above-50\% count (10 of 11 seasons), demonstrating the most consistent performance.

\section{Statistical Significance Testing}
With over 11,000 bets per model, even small deviations from $50\%$ require formal significance testing. We apply four
complementary tools: Wilson score confidence intervals, bootstrap resampling (standard and block bootstrap to
account for temporal correlation), Bonferroni-corrected $p$-values for multiple comparisons, and McNemar's paired test
for pairwise model comparison. \\
\begin{figure}[H]
    \centering
    \includegraphics[width=1\linewidth]{figure 6.png}
    \caption{Confidence intervals, significance testing (raw and Bonferroni-corrected), and per-season accuracy distributions.}
    \label{fig: 6}
\end{figure}

After Bonferroni correction for testing 6 models, four survive with significance against $50\%$: Ridge, Lasso, ElasticNet,
and PCR (all $p < 0.006$). XGBoost ($p = 0.054$) fails to reach significance even before correction. All three CI methods
(Wilson, standard bootstrap, block bootstrap) produce nearly identical intervals, confirming that temporal correlation
does not materially inflate confidence. However, every model's CI includes $52.4\%$, meaning we cannot confirm
profitability with flat betting. \\
\begin{figure}[H]
    \centering
    \includegraphics[width=1\linewidth]{figure 7.png}
    \caption{McNemar pairwise disagreement heatmap. Asterisks denote statistically significant differences ($p < 0.05$).}
    \label{fig: 7}
\end{figure}
McNemar's test confirms that Lasso significantly outperforms PCR, XGBoost, and RBF+Ridge in pairwise
comparisons. The linear models (Ridge, Lasso, ElasticNet) do not differ significantly from each other, consistent with
extracting similar signal from the features. \\

\section{Selective Betting with Confidence Thresholds}
Since no model is profitable when betting every game, we test whether restricting bets to games where the model
strongly disagrees with the Vegas line can push accuracy above break-even. The absolute difference between the
model's prediction and the Vegas total serves as a confidence measure. We sweep thresholds from $0.5$ to 20 points, subject to dual validity constraints: at least 300 total bets globally, and at least 10 bets in every individual season. \\
\begin{figure}[H]
    \centering
    \includegraphics[width=1\linewidth]{figure 8.png}
    \caption{Bet volume decay and ROI by confidence threshold for all six models. Stars mark optimal thresholds.}
    \label{fig: 8}
\end{figure}

Four models achieve positive ROI at valid thresholds. Ridge leads at $+5.4\%$ ROI ($55.0\%$ accuracy, 826 bets at
threshold $\geq 4.5$), followed by ElasticNet ($+2.2\%$, 1,896 bets at $\geq 4.0$), Lasso ($+2.1\%$, 4,529 bets at $\geq 2.0$), and PCR
($+0.6\%$, 5,126 bets at $\geq 2.5$). XGBoost and RBF+Ridge never achieve positive ROI at any valid threshold.\\
The explanation lies in prediction variance. Linear models have prediction standard deviations of 2.4-2.8 points,
making conservative adjustments to the efficient Vegas line. When they strongly disagree with Vegas, it likely reflects
genuine signal. XGBoost and RBF+Ridge have standard deviations of approximately 6 points, routinely making
extreme predictions driven by overfitting rather than genuine market mispricing. \\
\begin{figure}[H]
    \centering
    \includegraphics[width=1\linewidth]{figure 9.png}
    \caption{Comparison of base (all bets) vs. selective betting — volume, accuracy, and ROI for each model.}
    \label{fig: 9}
\end{figure}

After Bonferroni correction, five models remain significantly above $50\%$ at their optimal thresholds. However, no
model's selective accuracy is significantly above $52.4\%$. The evidence for profitability is suggestive but not
statistically conclusive, as the smaller sample sizes from selective betting widen the confidence intervals.

\section{Bankroll Simulation and Holdout Deployment}
The four qualifying models are subjected to bankroll simulations across 10 sizing rules, ranging from conservative flat
dollar amounts to the theoretically optimal Kelly criterion and its fractional variants. Starting from a \$10,000 bankroll,
the simulations run across all 11 walk-forward seasons to evaluate growth, drawdown, and consistency. \\
\begin{figure}[H]
    \centering
    \includegraphics[width=1\linewidth]{figure 10.png}
    \caption{Backtest heatmap of final bankroll values and CAGR across all model-strategy combinations.}
    \label{fig: 10}
\end{figure}

The simulations reveal that sizing matters as much as accuracy. Flat \$500 bets cause bankruptcy for all four models
despite positive edges. Full Kelly achieves the highest CAGR but with extreme drawdowns ($76.6\%$ for Ridge).
Fractional Kelly variants consistently offer better risk-adjusted returns. \\
\begin{figure}[H]
    \centering
    \includegraphics[width=1\linewidth]{figure 11.png}
    \caption{Risk-return scatter of all model-strategy combinations. The Pareto frontier runs through Ridge's conservative strategies.}
    \label{fig: 11}
\end{figure}

\subsection{Model Nomination}
Using only backtest data available before the 2024/25 season, we nominate Lasso with Flat 2\% sizing (\$200 per bet)
as the deployment strategy. While Ridge offers marginally higher CAGR ($+10.6\%$ vs. $+9.5\%$), Lasso wins on every
other criterion: lower maximum drawdown ($57.2\%$ vs. $76.6\%$), more winning seasons (7/11 vs. 6/11), $5.5$ times more
bets (4,529 vs. 826) for greater statistical reliability, and a simpler sizing rule that requires no win-probability
estimation. This nomination is entirely prospective.

\subsection{Holdout Results: 2024/25 Season}
The Lasso model places 256 qualifying bets on the unseen 2024/25 season, achieving $53.1\%$ accuracy. This is
within $0.3$ percentage points of its 11-season backtest accuracy of $53.4\%$, and sits above the $52.4\%$ break-even
threshold. Under Flat 2\% sizing, the \$10,000 bankroll grows to \$10,727, a $+7.3\%$ return with a maximum drawdown
of $15.3\%$ \\
\begin{figure}[H]
    \centering
    \includegraphics[width=1\linewidth]{figure 12.png}
    \caption{Lasso Flat 2\% deployment on the 2024/25 holdout season — bankroll trajectory and cumulative win rate convergence.}
    \label{fig: 12}
\end{figure}

The consistency between backtest and holdout is the most important finding: the strategy nominated prospectively on
the basis of 11 seasons of data remains profitable on genuinely unseen data. As an investment benchmark, the
strategy's $+7.3\%$ return compares to the S\&P 500's $+3.7\%$ over the same October 2024 to June 2025 window, with
near-zero correlation to equity markets. \\
A single season of 256 bets is statistically underpowered to prove a persistent edge. The 95\% Wilson confidence
interval for holdout accuracy spans [$47.0\%, 59.1\%$], which includes both 50\% and $52.4\%$. The positive result is
encouraging but not conclusive.

\section{Real-World Implementation Constraints}
Before drawing conclusions about practical viability, we examine three frictions that stand between a profitable
backtest and actual profits: execution slippage, market depth, and account restrictions.

\subsection{Execution Slippage}
We simulate scenarios where the betting line shifts $0.5$ to $3.0$ points against the bet direction before execution. At
zero slippage, Lasso achieves its full $+1.9\%$ flat-bet ROI on 4,529 backtest bets. The break-even slippage is
approximately $0.5$ points, which falls at the lower bound of typical NBA line movement ($0.5$-$1.5$ points). Even modest
execution delays can eliminate the advantage entirely. \\
\begin{figure}[H]
    \centering
    \includegraphics[width=1\linewidth]{figure 13.png}
    \caption{Accuracy and ROI degradation under increasing execution slippage. The typical slippage range (grey band) largely overlaps
with the unprofitable region.}
    \label{fig: 13}
\end{figure}

\subsection{Market Depth and Scalability}
Sportsbook betting limits cap the maximum bet per game at \$2,500-\$10,000 depending on the book. Across
approximately 10 accessible US sportsbooks, the maximum deployable capital per game is roughly \$50,000,
implying a ceiling bankroll of \$2.5 million at Flat 2\% sizing. A \$10,000 bankroll generates an expected annual profit of
approximately \$188; a \$250,000 bankroll generates approximately \$4,700. Scaling up capital does not fix the
slippage problem, as break-even slippage is a property of the edge, not the bankroll size. \\
\begin{figure}[H]
    \centering
    \includegraphics[width=1\linewidth]{figure 14.png}
    \caption{Expected profit by bankroll size and combined impact of slippage and capital deployment.}
    \label{fig: 14}
\end{figure}

\subsection{Account Restrictions}
The most significant threat is account restrictions. Sportsbooks track betting patterns and flag accounts that exhibit
Closing Line Value (consistently getting better odds than the final line). Winning accounts are typically limited or
closed within 1-3 seasons. Our Lasso model places approximately 412 bets per season at a thin $53.4\%$ edge, a
volume and pattern that would likely trigger restrictions within 2-3 seasons.

\section{Conclusion}
This project set out to determine whether machine learning models can profitably predict NBA over/under totals. The
answer is nuanced: a genuine statistical edge exists, but real-world frictions make it extremely difficult to capture
reliably.

\subsection{The Edge is Real}
Linear models (Ridge, Lasso, ElasticNet) achieve betting accuracies of $51.7$–$52.0\%$ on over 11,000 out-of-sample
predictions, with $p$-values below 0.001 after Bonferroni correction. Selective betting pushes four models above the
$52.4\%$ break-even threshold, with Ridge achieving $55.0\%$ accuracy and $+5.4\%$ ROI on 826 high-confidence bets.

\subsection{Simple Models Dominate}
Simple linear models consistently outperform complex nonlinear alternatives (XGBoost, RBF+Ridge). XGBoost
shows a 2.3-point train-vs-OOS RMSE gap, confirming that it memorizes training patterns that do not generalize. In
an efficient market where exploitable signal is small and noisy, the additional flexibility of complex models becomes a
liability rather than an asset. This is a textbook illustration of the bias-variance tradeoff in a low-signal environment.

\subsection{The Holdout Confirms the Backtest}
Our prospectively nominated Lasso Flat 2\% strategy achieved $53.1\%$ accuracy on 256 unseen 2024/25 bets, within
0.3 percentage points of its 11-season backtest accuracy of $53.4\%$. The \$10,000 bankroll grew to \$10,727 ($+7.3\%$
return), closely tracking the $+9.5\%$ annualized CAGR observed in the backtest. This consistency between
prospective nomination and holdout performance is the most important result of the project, though a single season
of 256 bets is insufficient for statistical certainty.

\subsection{Real-World Viability is Limited}
The strategy passes three of six viability criteria (backtest accuracy, backtest ROI, holdout confirmation) but fails the
remaining three. The break-even slippage of approximately $0.5$ points falls at the lower bound of typical NBA line
movement ($0.5$–$1.5$ points), meaning even modest execution delays can eliminate the edge. Sportsbook account
restrictions track betting patterns and flag accounts that exhibit Closing Line Value, imposing a hard time limit of 1–3
seasons before the account is limited or closed.

\subsection{Implications for Market Efficiency}
Our results are consistent with the semi-strong form of market efficiency: publicly available statistical information is
largely priced into the Vegas line. The residual edge ($+1.9\%$ flat-bet ROI in backtest, $+7.3\%$ bankroll ROI in the
holdout season) is real but small enough that transaction costs and market structure constraints make it extremely
difficult to capture reliably. From a machine learning perspective, this project demonstrates that model accuracy and
economic viability are fundamentally different questions. A model can be statistically superior to random chance ($p <0.001$) yet still fail to generate profit after accounting for the market's built-in margin. The gap between statistical
significance and economic significance is the defining challenge of applied ML in efficient markets.

\end{document}
